% Self-contained source: no external .bib file or figures are required.
% Compile twice with: pdflatex "kiem_lee_e_positivity(1).tex"
\documentclass[11pt]{article}
\usepackage[T1]{fontenc}
\usepackage{lmodern}
\usepackage[margin=1in]{geometry}
\usepackage{amsmath,amssymb,amsthm,mathtools}
\usepackage{microtype}
\usepackage[hidelinks]{hyperref}
\hypersetup{
  pdftitle={Elementary positivity for vertex multicolorings and dual Dyck functions},
  pdfsubject={Weighted chromatic positivity and its applications to generalized Hessenberg and step-tau dual Dyck symmetric functions},
  pdfauthor={}
}
\newtheorem{theorem}{Theorem}[section]
\newtheorem{lemma}[theorem]{Lemma}
\newtheorem{corollary}[theorem]{Corollary}
\theoremstyle{remark}
\newtheorem{remark}[theorem]{Remark}
\newcommand{\Z}{\mathbb Z}
\newcommand{\NN}{\mathbb Z_{\geq 0}}
\newcommand{\Npos}{\mathbb Z_{>0}}
\newcommand{\xx}{\mathbf x}
\newcommand{\avec}{\mathbf a}
\newcommand{\bvec}{\mathbf b}
\newcommand{\cvec}{\mathbf c}
\newcommand{\one}{\mathbf 1}
\newcommand{\cA}{\mathcal A}
\DeclareMathOperator{\asc}{asc}
\DeclareMathOperator{\csf}{csf}
\DeclareMathOperator{\di}{di}
\DeclareMathOperator{\dinv}{dinv}
\DeclareMathOperator{\DS}{DS}
\DeclareMathOperator{\supp}{supp}
\numberwithin{equation}{section}
\setlength{\emergencystretch}{1.5em}
\setlength{\parskip}{2pt}
\setlength{\parindent}{1.2em}

\begin{document}
\begin{center}
  {\Large\bfseries Elementary positivity for vertex multicolorings\\
  and dual Dyck functions}\par
\end{center}

\begin{abstract}
We deduce coefficientwise elementary positivity for vertex multicolorings
of natural unit interval graphs from the intermediate monomial-positivity
results in OpenAI's chromatic-positivity preprint. A vertex of multiplicity
$a_i$ receives a set of $a_i$ distinct colors. Extracting arbitrary
nonnegative vertex exponents, rather than only squarefree exponents,
extends the chromatic argument to these multiplicities without division
by quantum factorials. This yields positivity for generalized Hessenberg
chromatic functions. In a second part, a direct correspondence between
step-$\tau$ dual Dyck factorizations and multicolorings proves elementary
positivity in each statistic degree for every integer $\tau\geq0$.
An explicit grading shift recovers Hawkes's original dual Dyck symmetric
functions at step one. All positivity conclusions use the stated external
input.
\end{abstract}

\section{Introduction}
Elementary positivity asks for a symmetric function to be a nonnegative
linear combination of products of elementary symmetric functions. For
chromatic functions refined by an ascent statistic, the coefficientwise
question requires the elementary coefficients to be polynomials with
nonnegative integer coefficients. We consider the multicoloring version:
each vertex receives a prescribed number of distinct colors, and adjacent
vertices receive disjoint color sets. This construction appears in
Hwang~\cite{Hwang} and Kiem--Lee~\cite{KL}; the multiplicities specify
color-set cardinalities rather than powers of a single color.

The first part proves elementary positivity for every multiplicity vector
on a natural unit interval graph (Theorem~\ref{thm:main}). The central
observation is that the monomial-positivity results of~\cite{OAI} apply to
all nonnegative vertex exponents. Restriction from a triangular graph and
the elementary--monomial Cauchy identity then give the weighted theorem.
This is a deduction from the specified intermediate results of~\cite{OAI},
not an independent proof of chromatic positivity. The generalized
Hessenberg functions of Kiem--Lee are one application, giving an
affirmative answer to their Question~4.9 under the same input
(Corollary~\ref{cor:KL}).

The second part extends the dual Dyck factor conditions of
Hawkes~\cite{Hawkes} to every integer step $\tau\geq0$. These
factorizations encode multicolorings of the graph joining distinct
integer values at distance at most $\tau$. Repeated values become vertex
multiplicities, and the statistic is a constant shift of the complement
of the weighted ascent statistic. Theorem~\ref{thm:dyck} gives elementary
positivity in each statistic degree. At step one, removing the
equal-pair shift recovers Hawkes's original functions
(Corollary~\ref{cor:classical}), strengthening their Schur positivity
in~\cite[Corollary~3.19]{Hawkes}.

The multiplicity step cannot be obtained simply by dividing an unweighted
positive expansion. Replacing each vertex by a clique of its multiplicity
introduces a factor $\prod_i[a_i]_q!$, where
$[a]_q!=\prod_{j=1}^a(1+q+\cdots+q^{j-1})$; see
\cite[Remark~4.4]{KL}. Division by such a factor need not preserve
coefficientwise positivity, as $(1+q)(1-q+q^2)=1+q^3$ illustrates.
The coefficient extraction below counts color sets directly.

\section{Definitions and conventions}\label{sec:definitions}
Let $\Lambda_{\Z}=\Z[e_1,e_2,\ldots]$ be the ring of symmetric functions in
$\xx=(x_1,x_2,\ldots)$, where
$e_k(\xx)=\sum_{c_1<\cdots<c_k}x_{c_1}\cdots x_{c_k}$ and $e_0=1$.
For a partition $\lambda$, write $|\lambda|$ for its size,
$\ell(\lambda)$ for its length, $e_\lambda=\prod_j e_{\lambda_j}$,
and $m_\lambda$ for the sum of distinct monomials with positive exponent
multiset $\lambda$. A homogeneous symmetric function is
\emph{$e$-positive over $\NN[q]$} when its $e_\lambda$-coefficients lie in
$\NN[q]$.

An ordered graph $G$ on $[r]=\{1,\ldots,r\}$ is a \emph{natural unit
interval graph} if, for a weakly increasing map $g:[r]\to[r]$ with
$g(i)\geq i$, its edges are $\{i,j\}$ with $i<j\leq g(i)$.
For $\avec\in\NN^r$, set
\begin{equation}\label{eq:NM}
 N=|\avec|:=\sum_i a_i,
 \qquad M=M_G(\avec):=\sum_{\substack{i<j\\\{i,j\}\in E(G)}}a_i a_j.
\end{equation}
A \emph{proper multicoloring of multiplicity $\avec$} assigns a finite set
$C_i\subseteq\Npos$ to each vertex, with $|C_i|=a_i$ and
$C_i\cap C_j=\varnothing$ on every edge. Define
\begin{equation}\label{eq:weighted-function}
 \begin{aligned}
 \asc_G(C)&=\sum_{\substack{i<j\\\{i,j\}\in E(G)}}
       \#\{(c,d)\in C_i\times C_j:c<d\},\\
 X_{G,\avec}(\xx;q)&=\sum_{C\text{ proper}}q^{\asc_G(C)}
       \prod_{i=1}^r\prod_{c\in C_i}x_c.
 \end{aligned}
\end{equation}

Put $R=\Z[v,v^{-1}]$. The \emph{nonnegative-exponent graph algebra} is
$\cA_G^+=\bigoplus_{\bvec\in\NN^r} RZ^{\bvec}$, with multiplication
\begin{equation}\label{eq:algebra}
 Z^{\bvec}Z^{\cvec}=v^{\Omega_G(\bvec,\cvec)}Z^{\bvec+\cvec},
 \qquad
 \Omega_G(\bvec,\cvec)=\sum_{\substack{i<j\\\{i,j\}\in E(G)}}
                  (b_jc_i-b_ic_j).
\end{equation}
Bilinearity of $\Omega_G$ gives associativity. The $Z^{\bvec}$ are basis
symbols, and $[Z^{\bvec}]$ denotes coefficient extraction in this basis.
For an independent set $J$, let $\one_J$ be its indicator vector and put
\begin{equation}\label{eq:independent-sums}
 E_k(G)=\sum_{\substack{J\subseteq[r]\text{ independent}\\|J|=k}}
           Z^{\one_J},\qquad E_0(G)=1.
\end{equation}
By \cite[Proposition~2.2]{OAI}, these elements commute. Thus
$\Phi_G:\Lambda_{\Z}\to\cA_G^+$ is the homomorphism determined by
$\Phi_G(e_k)=E_k(G)$. Its degree is the total vertex exponent.
All graph notation also applies to any finite totally ordered vertex set
by increasing relabeling.

\section{Elementary positivity for vertex multicolorings}\label{sec:proof}
\noindent\textit{External dependency.}
We take Proposition~2.2, Lemmas~5.1 and~5.7, and Theorem~5.8 of the
September~24, 2026 version of \cite{OAI} as external results.
The proof below is a deduction from those assertions; it does not
independently establish their validity.

\begin{theorem}\label{thm:main}
Assuming the cited external results, for every natural unit interval
graph $G$ and every $\avec\in\NN^{V(G)}$,
\begin{equation}\label{eq:main}
 X_{G,\avec}(\xx;q)=\sum_{\lambda\vdash N}
                  c_\lambda(G,\avec;q)e_\lambda(\xx),
 \qquad c_\lambda(G,\avec;q)\in\NN[q].
\end{equation}
More precisely, with $M$ as in \eqref{eq:NM},
\begin{equation}\label{eq:coefficient-formula}
 c_\lambda(G,\avec;q)
   =\sum_{d=0}^{M}
       \bigl[v^{\,2d-M}Z^{\avec}\bigr]\Phi_G(m_\lambda)\,q^d.
\end{equation}
In particular, each $c_\lambda(G,\avec;q)$ has degree at most $M$.
\end{theorem}

\subsection{Positivity at arbitrary vertex multiplicities}
\begin{lemma}\label{lem:positivity}
Under the external results above, for every natural unit interval graph
$G$, every partition $\lambda$, and every $\avec\in\NN^{V(G)}$,
\begin{equation}\label{eq:positive-coefficient}
 [Z^{\avec}]\Phi_G(m_\lambda)\in\NN[v,v^{-1}].
\end{equation}
\end{lemma}
\begin{proof}
Choose an ordered induced embedding of $G$ into an auxiliary triangular
graph $\Gamma$ supplied by \cite[Lemma~5.1]{OAI}. Write $D$ for its
independence number. The two results
\cite[Lemma~5.7 and Theorem~5.8]{OAI}, applied in the all-negative chamber
of that construction, give
\begin{equation}\label{eq:ambient-positivity}
 \Phi_\Gamma(m_\lambda)
       \in\bigoplus_{\bvec\in\NN^{V(\Gamma)}}
                       \NN[v,v^{-1}]Z^{\bvec}
       \qquad\text{if }\ell(\lambda)\leq D.
\end{equation}
Specifically, Theorem~5.8 identifies this evaluation with the section
whose finiteness, nonnegative vertex support, and Laurent-coefficient
positivity are asserted in Lemma~5.7. Thus \eqref{eq:ambient-positivity}
concerns every vertex exponent, not only exponents with coordinates
zero or one.

To compare their finite-alphabet convention with ours, note that
$E_k(\Gamma)=0$ for $k>D$. Hence $\Phi_\Gamma$ factors through the
specialization of symmetric functions to $D$ variables. This agrees with
the evaluation in the cited results. If $\ell(\lambda)>D$, that
specialization sends $m_\lambda$ to zero, so
\eqref{eq:ambient-positivity} holds for these partitions as well.

Identify $V(G)$ with the selected vertices of $\Gamma$ and define the
$R$-linear map
\begin{equation}\label{eq:projection}
 \rho:\cA_\Gamma^+\longrightarrow\cA_G^+,
 \qquad
 \rho(Z^{\bvec})=
 \begin{cases}
  Z^{\bvec|_{V(G)}},&\supp(\bvec)\subseteq V(G),\\
  0,&\text{otherwise},
 \end{cases}
\end{equation}
where $\supp(\bvec)=\{u:b_u\ne0\}$.
This is an algebra homomorphism. Indeed, if either of two nonnegative
exponent vectors uses an outside vertex, their sum still uses that
vertex; both sides of multiplicativity are then zero. If both are
supported on $V(G)$, the induced embedding gives
$\Omega_\Gamma(\bvec,\cvec)=
 \Omega_G(\bvec|_{V(G)},\cvec|_{V(G)})$, which proves multiplicativity
in the remaining case.

An independent set of $\Gamma$ contained in $V(G)$ is exactly an
independent set of $G$. Therefore
$\rho(E_k(\Gamma))=E_k(G)$ for every $k$, and equality on the
algebraically independent generators $e_k$ implies
\begin{equation}\label{eq:intertwining}
 \rho\circ\Phi_\Gamma=\Phi_G.
\end{equation}
Let $\widetilde{\avec}$ extend $\avec$ by zero outside $V(G)$. Then
\[
 [Z^{\avec}]\Phi_G(m_\lambda)
   =[Z^{\widetilde{\avec}}]\Phi_\Gamma(m_\lambda)
   \in\NN[v,v^{-1}],
\]
as required. Crucially, \eqref{eq:projection} is defined on the
nonnegative-exponent algebra, not on the full quantum torus, where
vertex generators would be invertible.
\end{proof}

\subsection{The multicoloring kernel and its normalization}
For $s\geq1$, write $\xx_s=(x_1,\ldots,x_s)$ and set
\begin{equation}\label{eq:kernel}
 \mathcal E_G(t)=\sum_{k\geq0}E_k(G)t^k,
 \qquad
 K_s=\mathcal E_G(x_s)\mathcal E_G(x_{s-1})\cdots
                       \mathcal E_G(x_1).
\end{equation}
The sum defining $\mathcal E_G$ is finite. Its values commute, but the
displayed descending color order will match the ascent convention.

\begin{lemma}\label{lem:kernel}
For every $s\geq1$,
\begin{equation}\label{eq:kernel-identity}
 X_{G,\avec}(\xx_s;v^2)=v^M[Z^{\avec}]K_s.
\end{equation}
\end{lemma}
\begin{proof}
Expanding $K_s$ chooses an independent set $J_c$ for each color
$c\in[s]$. The associated term is
\[
 \left(\prod_{c=1}^s x_c^{|J_c|}\right)
 Z^{\one_{J_s}}Z^{\one_{J_{s-1}}}\cdots Z^{\one_{J_1}}.
\]
Its vertex exponent is $\sum_c\one_{J_c}$. Thus extraction of
$Z^{\avec}$ is equivalent to requiring
$C_i:=\{c\in[s]:i\in J_c\}$ to have cardinality $a_i$.
Independence of each $J_c$ says exactly that adjacent vertices receive
disjoint color sets. Conversely, the sets $C_i$ recover all the $J_c$
uniquely. This is a bijection with the multicolorings in
\eqref{eq:weighted-function}; no occurrences at a vertex are labeled.

Repeated use of \eqref{eq:algebra} gives the Laurent weight
$v^{T(C)}$, where
\begin{equation}\label{eq:energy}
 T(C)=\sum_{c>d}\Omega_G(\one_{J_c},\one_{J_d}).
\end{equation}
Fix an edge $\{i,j\}$ with $i<j$ and a pair
$(a,b)\in C_i\times C_j$. The colors are unequal. If $a<b$, the
factor containing vertex $j$ occurs to the left of the factor
containing $i$, and its contribution to \eqref{eq:energy} is $+1$.
If $a>b$, the contribution is $-1$. There are $a_i a_j$ such pairs
on this edge. Consequently,
\begin{equation}\label{eq:normalization}
 T(C)=\asc_G(C)-\bigl(M-\asc_G(C)\bigr)
     =2\asc_G(C)-M.
\end{equation}
Multiplication by $v^M$ changes the weight to
$v^{2\asc_G(C)}$, proving \eqref{eq:kernel-identity}.
\end{proof}

\subsection{Extraction of the elementary coefficients}
\begin{proof}[Proof of Theorem~\ref{thm:main}]
If $N=0$, the unique multicoloring has all color sets empty and the
function is $1$; assume $N>0$. In a second commuting alphabet
$\mathbf y$, the elementary--monomial Cauchy identity is
\begin{equation}\label{eq:cauchy}
 \prod_{c=1}^s\left(\sum_{k\geq0}e_k(\mathbf y)x_c^k\right)
   =\prod_{c,j}(1+x_cy_j)
   =\sum_\lambda e_\lambda(\xx_s)m_\lambda(\mathbf y).
\end{equation}
For clarity, a fixed monomial in the $y_j$ with positive exponents
$\lambda_1,\ldots,\lambda_\ell$ has coefficient
$\prod_{i=1}^{\ell}e_{\lambda_i}(\xx_s)$ in the middle product.
Grouping such monomials proves the last equality. The identity is
interpreted degree by degree, so there is no convergence issue.

Apply $\Phi_G$ to the $\mathbf y$-symmetric functions. Since the
$E_k(G)$ commute, their generating factors may be put in the descending
order of \eqref{eq:kernel}, giving
\begin{equation}\label{eq:e-kernel}
 K_s=\sum_\lambda e_\lambda(\xx_s)\Phi_G(m_\lambda).
\end{equation}
The evaluation preserves total degree. Extracting $Z^{\avec}$ therefore
retains only partitions of $N$. Combining \eqref{eq:e-kernel} with
Lemma~\ref{lem:kernel}, we obtain, for every $s$,
\begin{equation}\label{eq:laurent-expansion}
 X_{G,\avec}(\xx_s;v^2)
   =\sum_{\lambda\vdash N}P_\lambda(v)e_\lambda(\xx_s),
 \qquad
 P_\lambda(v)=v^M[Z^{\avec}]\Phi_G(m_\lambda).
\end{equation}
Lemma~\ref{lem:positivity} gives
$P_\lambda(v)\in\NN[v,v^{-1}]$.
It remains to justify that these Laurent polynomials use only the even,
nonnegative powers $1,v^2,\ldots,v^{2M}$; Laurent positivity alone
would not establish the required conclusion over $\NN[q]$.

Take $s\geq N$. The polynomials
$\{e_\lambda(\xx_s):\lambda\vdash N\}$ form a basis of the homogeneous
symmetric polynomials of degree $N$. Every multicoloring satisfies
$0\leq\asc_G(C)\leq M$. Hence the left side of
\eqref{eq:laurent-expansion} has no power $v^j$ with
$j\notin\{0,2,\ldots,2M\}$. Extracting such a power gives
\[
 0=\sum_{\lambda\vdash N}
           [v^j]P_\lambda(v)\,e_\lambda(\xx_s).
\]
Linear independence implies $[v^j]P_\lambda(v)=0$ for every
$\lambda$. Thus there is a polynomial $c_\lambda(q)\in\NN[q]$
of degree at most $M$ such that $P_\lambda(v)=c_\lambda(v^2)$.
Its coefficients are exactly those in
\eqref{eq:coefficient-formula}. The substitution $q\mapsto v^2$
is injective, so \eqref{eq:laurent-expansion} proves
\eqref{eq:main} in finitely many colors. These coefficients are
independent of $s$, and setting $x_s=0$ is compatible with the
multicoloring definition. Passing to symmetric functions proves
\eqref{eq:main} in the infinite alphabet.
\end{proof}

\begin{remark}[Why no factorial division is needed]
At each vertex, the proof chooses a set of colors, not an ordering of
those colors. Hence the coefficient $[Z^{\avec}]$ already counts exactly
the required multicolorings. The essential input is the all-exponent
positivity of Lemma~\ref{lem:positivity}.
\end{remark}

\subsection{Generalized Hessenberg chromatic functions}
Take $0=i_0<i_1<\cdots<i_r=n$ and
$I=\{i_1,\ldots,i_{r-1}\}$. A \emph{generalized Hessenberg function}
is a weakly increasing map $h:I\cup\{n\}\to I\cup\{n\}$ with
$h(i)\geq i$. Its ordered weighted graph has vertices $I\cup\{n\}$,
edges $\{i,j\}$ for $i<j\leq h(i)$, and weight
$w_h(i_k)=i_k-i_{k-1}$. Kiem--Lee's function
$\csf_q(h)$ is the multicoloring series \eqref{eq:weighted-function}
for this graph and weight \cite[Definitions~2.1 and~4.3]{KL}.

\begin{corollary}\label{cor:KL}
Under the external input of Theorem~\ref{thm:main}, $\csf_q(h)$ is
$e$-positive over $\NN[q]$ for every generalized Hessenberg function $h$.
\end{corollary}
\begin{proof}
Define $g:[r]\to[r]$ by the requirement

\begin{equation}\label{eq:relabel}
 h(i_k)=i_{g(k)},\qquad a_k=i_k-i_{k-1}.
\end{equation}
Since $h$ is extensive and weakly increasing, $g(k)\geq k$ and
$g$ is weakly increasing. Increasing relabeling $i_k\mapsto k$
identifies the graph of $h$ with the natural unit interval graph $G_g$:
for $k<\ell$,
\[
 i_\ell\leq h(i_k)
 \quad\Longleftrightarrow\quad
 \ell\leq g(k).
\]
The same relabeling preserves the color-set cardinalities and every
ascending color pair across every edge. Accordingly,
\begin{equation}\label{eq:KL-identification}
 \csf_q(h)=X_{G_g,\avec}(\xx;q),\qquad
 |\avec|=\sum_{k=1}^r(i_k-i_{k-1})=n.
\end{equation}
Applying \eqref{eq:main} to \eqref{eq:KL-identification} yields
\[
 \csf_q(h)=\sum_{\lambda\vdash n}
       c_\lambda(G_g,\avec;q)e_\lambda(\xx),
 \qquad c_\lambda(G_g,\avec;q)\in\NN[q].
\]
This is precisely the elementary-positivity assertion of
\cite[Question~4.9]{KL}.
\end{proof}

\section{Elementary positivity for step-\texorpdfstring{$\tau$}{tau} dual Dyck functions}\label{sec:dyck}

Fix an integer $\tau\geq0$. A \emph{dual step-$\tau$ word} is a finite
integer word $w$ satisfying $w_{j+1}>w_j+\tau$ at each successive
position. Empty and one-letter words are allowed. For a finite multiset
$S$ of integers, a dual step-$\tau$ factorization is a sequence
$\mathcal F=(F_1,F_2,\ldots)$ of such words, all but finitely many empty,
whose concatenation $w(\mathcal F)$ rearranges $S$. Equal occurrences
are not labeled. The step-one factor conditions are those of
\cite[Definitions~2.3--2.4]{Hawkes}, with factor and variable indices
shifted from $0,1,\ldots$ to $1,2,\ldots$.

For integers $a,b$, define
\begin{equation}\label{eq:pair-stat}
 d_\tau(a,b)=
 \begin{cases}
  \max(0,a+\tau-b),&a\leq b,\\
  \max(0,b+1+\tau-a),&a>b,
 \end{cases}
 \qquad
 \dinv_\tau(w)=\sum_{p<q}d_\tau(w_p,w_q).
\end{equation}
Define the factorization series and its fixed-degree coefficients by
\begin{equation}\label{eq:dual-dyck}
 \begin{aligned}
 \DS^{*(\tau)}(S,d;\xx)
   &=\sum_{\substack{\mathcal F\text{ dual step-}\tau\\
                     \dinv_\tau(w(\mathcal F))=d}}
                   \prod_{a\geq1}x_a^{|F_a|},\\
 \mathcal D^*_{S,\tau}(\xx;q)
   &=\sum_{d\geq0}q^d\DS^{*(\tau)}(S,d;\xx).
 \end{aligned}
\end{equation}
Each monomial coefficient is finite because its exponents prescribe all
factor lengths. There is no condition on the first entry or sign of a
factor, and no restriction on $S$ beyond finiteness.

Write a nonempty multiset as $S=c_1^{a_1}\cdots c_r^{a_r}$, with
$c_1<\cdots<c_r$ and $a_i>0$. Let $G_{S,\tau}$ have vertices $[r]$
and edges
\begin{equation}\label{eq:dyck-graph}
 \{i,j\}\in E(G_{S,\tau})\quad\Longleftrightarrow\quad
 i<j\text{ and }c_j-c_i\leq\tau.
\end{equation}
This is a natural unit interval graph: its defining function is
$g(i)=\max\{j:c_j\leq c_i+\tau\}$, which is weakly increasing and
satisfies $g(i)\geq i$. Put
\begin{equation}\label{eq:dyck-shift}
 \begin{aligned}
 M_\tau(S)&=\sum_{\substack{i<j\\c_j-c_i\leq\tau}}a_i a_j,\\
 \beta_\tau(S)&=\tau\sum_i\binom{a_i}{2}
   +\sum_{\substack{i<j\\c_j-c_i\leq\tau}}
          (\tau-c_j+c_i)a_i a_j.
 \end{aligned}
\end{equation}
Both quantities are nonnegative integers.

\begin{lemma}\label{lem:dyck-statistic}
For every rearrangement $w$ of $S$,
\begin{equation}\label{eq:static-statistic}
 \dinv_\tau(w)=\beta_\tau(S)
       +\#\{(p,q):p<q,\ 0<w_p-w_q\leq\tau\}.
\end{equation}
\end{lemma}
\begin{proof}
Each equal pair contributes $\tau$. For distinct values with difference
$\delta>0$, the increasing order contributes $\max(0,\tau-\delta)$
and the decreasing order contributes $\max(0,\tau+1-\delta)$.
Since $\delta$ and $\tau$ are integers, both are zero when
$\delta>\tau$; otherwise they are $\tau-\delta$ and
$\tau-\delta+1$. Summing the order-independent contributions gives
$\beta_\tau(S)$, with one extra unit for each counted decreasing pair.
\end{proof}

\begin{theorem}\label{thm:dyck}
Under the external input of Theorem~\ref{thm:main}, for every finite
multiset $S$ of integers and all integers $\tau,d\geq0$, the series
$\DS^{*(\tau)}(S,d;\xx)$ is symmetric and $e$-positive over $\NN$.
More precisely, for nonempty $S$, write $G=G_{S,\tau}$,
$\avec=(a_1,\ldots,a_r)$, $M=M_\tau(S)$, and $\beta=\beta_\tau(S)$.
Then
\begin{equation}\label{eq:dyck-chromatic}
 \begin{aligned}
 \mathcal D^*_{S,\tau}(\xx;q)
   &=q^{\beta+M}X_{G,\avec}(\xx;q^{-1})\\
   &=\sum_{\lambda\vdash |S|}
       q^{\beta+M}c_\lambda(G,\avec;q^{-1})e_\lambda(\xx),
 \end{aligned}
\end{equation}
and each displayed elementary coefficient belongs to $\NN[q]$.
In particular, for $\beta\leq d\leq\beta+M$,
\begin{equation}\label{eq:dyck-coefficient}
 [e_\lambda]\DS^{*(\tau)}(S,d;\xx)
    =[q^{\beta+M-d}]c_\lambda(G,\avec;q)\in\NN,
\end{equation}
and $\DS^{*(\tau)}(S,d;\xx)=0$ outside that interval.
\end{theorem}
\begin{proof}
Every dual step-$\tau$ factor lists distinct values in increasing order
with successive gaps greater than $\tau$. All its pairwise gaps
therefore exceed $\tau$, so its value levels form an independent set of $G$.
Assign to vertex $i$ the set
$C_i=\{a:c_i\text{ occurs in }F_a\}$. This gives a proper
multicoloring with $|C_i|=a_i$. Conversely, from such a multicoloring,
form $F_a$ by listing in increasing order all $c_i$ with $a\in C_i$.
Independence ensures the required gaps. These constructions are inverse
and preserve monomials, including empty factors and repeated values.

For an edge $i<j$, occurrences in factors $a\in C_i$ and $b\in C_j$
form a decreasing pair in~\eqref{eq:static-statistic} exactly when
$b<a$. Properness excludes $a=b$, so among the $a_i a_j$ pairs on
this edge, these are exactly those not counted by $\asc_G$. Every
decreasing pair counted there comes from an edge. Consequently,
Lemma~\ref{lem:dyck-statistic} gives
\[
 \dinv_\tau(w(\mathcal F))=\beta+M-\asc_G(C).
\]
Summing this monomial-preserving identity proves the first equality
in~\eqref{eq:dyck-chromatic}. By Theorem~\ref{thm:main}, each
$c_\lambda(G,\avec;q)$ has nonnegative integer coefficients and
degree at most $M$. Reversing its coefficients and multiplying by
$q^\beta$ therefore gives a positive polynomial supported in
$[\beta,\beta+M]$. Extracting $q^d$ proves the claims. For $S=\varnothing$,
the unique factorization has every factor empty, so
$\DS^{*(\tau)}(\varnothing,0;\xx)=1$ and all other degrees are zero.
\end{proof}

\subsection{The original dual Dyck functions}
Hawkes's original statistic is
\[
 \di(w)=\#\{(p,q):p<q,\ w_p=w_q+1\};
\]
in particular, equal pairs are not counted
\cite[Definition~2.1]{Hawkes}. Put $\kappa(S)=\sum_i\binom{a_i}{2}$.
At $\tau=1$, the distinct-value terms in $\beta_1(S)$ vanish and
\eqref{eq:static-statistic} becomes
\begin{equation}\label{eq:classical-shift}
 \beta_1(S)=\kappa(S),\qquad
 \dinv_1(w)=\kappa(S)+\di(w).
\end{equation}

\begin{corollary}\label{cor:classical}
Under the same external input, the original dual Dyck symmetric functions
$\DS^*(S,d;\xx)$ of~\cite[Definition~2.5]{Hawkes} are $e$-positive
over $\NN$ for every finite multiset $S$ and every $d\geq0$.
\end{corollary}
\begin{proof}
The step-one factor conditions agree with Hawkes's, after reindexing
the variables. Equation~\eqref{eq:classical-shift} gives
\[
 \DS^*(S,d;\xx)=\DS^{*(1)}(S,d+\kappa(S);\xx),
\]
so Theorem~\ref{thm:dyck} applies. The empty multiset is immediate.
\end{proof}

For $\tau=0$, the value graph has no edges and both $M_0(S)$ and
$\beta_0(S)$ vanish, giving $\mathcal D^*_{S,0}=\prod_i e_{a_i}$.
For a single value repeated $m$ times, the graph has no edges at any
step, but the equal-pair shift gives
$\mathcal D^*_{S,\tau}=q^{\tau\binom m2}e_m$.
For $\tau>1$, nearby distinct values can also contribute to
$\beta_\tau(S)$; the equal-pair term alone would not give the correct
grading.

\begin{thebibliography}{OAI}

\bibitem[Haw]{Hawkes}
Graham Hawkes,
\emph{Dyck symmetric functions and applications to $q,t$-Catalan polynomials},
preprint, 2026.
\href{https://arxiv.org/abs/2605.13003v1}{arXiv:2605.13003v1}.
Numbered references above follow this version.

\bibitem[Hwa]{Hwang}
Byung-Hak Hwang,
\emph{Chromatic quasisymmetric functions and noncommutative
$P$-symmetric functions},
Transactions of the American Mathematical Society \textbf{377} (2024),
2855--2896.
\href{https://doi.org/10.1090/tran/9096}{doi:10.1090/tran/9096}.
\href{https://arxiv.org/abs/2208.09857v2}{arXiv:2208.09857v2}.

\bibitem[KL]{KL}
Young-Hoon Kiem and Donggun Lee,
\emph{Birational geometry of generalized Hessenberg varieties and the
generalized Shareshian--Wachs conjecture},
Journal of Combinatorial Theory, Series A \textbf{206} (2024), 105884.
\href{https://doi.org/10.1016/j.jcta.2024.105884}{doi:10.1016/j.jcta.2024.105884}.
\href{https://arxiv.org/abs/2208.12282v3}{arXiv:2208.12282v3},
19 February 2024. Numbered references above follow this arXiv version.

\bibitem[OAI]{OAI}
OpenAI,
\emph{Elementary positivity of chromatic quasisymmetric functions},
OpenAI Math Release preprint, 24 September 2026, 51 pp.
\href{https://github.com/openai/math/blob/fd4aeeb/preprints/Elementary-Positivity-of-Chromatic-Quasisymmetric-Functions-September-24-2026/paper.pdf}{OpenAI mathematics repository: manuscript PDF at revision \texttt{fd4aeeb}}.
The external inputs used here are Proposition~2.2, Lemmas~5.1 and~5.7,
and Theorem~5.8 of that version.

\end{thebibliography}
\end{document}
